\chapter{The Abbé Fouquet}

\lettrine{F}{ouquet} hastened back to his apartment by the subterranean passage, and immediately closed the mirror with the spring. He was scarcely in his well-known voice crying:—<Open the door, monseigneur, I entreat you, open the door!> Fouquet quickly restored a little order to everything that might have revealed either his absence or his agitation: he spread his papers over the desk, took up a pen, and, to gain time, said, through the closed door,—<Who is there?>

<What, monseigneur, do you not know me?> replied the voice.

<Yes, yes,> said Fouquet to himself, <yes, my friend, I know you well enough.> And then, aloud: <Is it not Gourville?>

<Why, yes, monseigneur.>

Fouquet arose, cast a look at one of his glasses, went to the door, pushed back the bolt, and Gourville entered. <Ah! monseigneur! monseigneur!> cried he, <what cruelty!>

<In what?>

<I have been a quarter of an hour imploring you to open the door, and you would not even answer me.>

<Once and for all, you know that I will not be disturbed when I am busy. Now, although I might make you an exception, Gourville, I insist upon my orders being respected by others.>

<Monseigneur, at this moment, orders, doors, bolts, locks, and walls I could have broken, forced and overthrown!>

<Ah! ah! it relates to some great event, then?> asked Fouquet.

<Oh! I assure you it does, monseigneur,> replied Gourville.

<And what is this event?> said Fouquet, a little troubled by the evident agitation of his most intimate confidant.

<There is a secret chamber of justice instituted, monseigneur.>

<I know there is, but do the members meet, Gourville?>

<They not only meet, but they have passed a sentence, monseigneur.>

<A sentence?> said the superintendent, with a shudder and pallor he could not conceal. <A sentence!—and on whom?>

<Two of your best friends.>

<Lyodot and D'Eymeris, do you mean? But what sort of a sentence?>

<Sentence of death.>

<Passed? Oh! you must be mistaken, Gourville; that is impossible.>

<Here is a copy of the sentence which the king is to sign to-day, if he has not already signed it.>

Fouquet seized the paper eagerly, read it, and returned it to Gourville. <The king will never sign that,> said he.

Gourville shook his head.

<Monseigneur, M.~Colbert is a bold councillor: do not be too confident!>

<Monsieur Colbert again!> cried Fouquet. <How is it that that name rises upon all occasions to torment my ears, during the last two or three days? You make so trifling a subject of too much importance, Gourville. Let M.~Colbert appear, I will face him; let him raise his head, I will crush him; but you understand, there must be an outline upon which my look may fall, there must be a surface upon which my feet may be placed.>

<Patience, monseigneur; for you do not know what Colbert is—study him quickly; it is with this dark financier as it is with meteors, which the eye never sees completely before their disastrous invasion; when we feel them we are dead.>

<Oh! Gourville, this is going too far,> replied Fouquet, smiling; <allow me, my friend, not to be so easily frightened; M.~Colbert a meteor! \swear{Corbleu,} we confront the meteor. Let us see acts, and not words. What has he done?>

<He has ordered two gibbets of the executioner of Paris,> answered Gourville.

Fouquet raised his head, and a flash gleamed from his eyes. <Are you sure of what you say?> cried he.

<Here is the proof, monseigneur.> And Gourville held out to the superintendent a note communicated by a certain secretary of the Hôtel de Ville, who was one of Fouquet's creatures.

<Yes, that is true,> murmured the minister; <the scaffold may be prepared, but the king has not signed; Gourville, the king will not sign.>

<I shall soon know,> said Gourville.

<How?>

<If the king has signed, the gibbets will be sent this evening to the Hôtel de Ville, in order to be got up and ready by to-morrow morning.>

<Oh! no, no!> cried the superintendent, once again; <you are all deceived, and deceive me in my turn; Lyodot came to see me only the day before yesterday; only three days ago I received a present of some Syracuse wine from poor D'Eymeris.>

<What does that prove?> replied Gourville, <except that the chamber of justice has been secretly assembled, has deliberated in the absence of the accused, and that the whole proceeding was complete when they were arrested.>

<What! are they, then, arrested?>

<No doubt they are.>

<But where, when, and how have they been arrested?>

<Lyodot, yesterday at daybreak; D'Eymeris, the day before yesterday, in the evening, as he was returning from the house of his mistress; their disappearances had disturbed nobody; but at length M.~Colbert all at once raised the mask, and caused the affair to be published; it is being cried by sound of trumpet, at this moment in Paris, and, in truth, monseigneur, there is scarcely anybody but yourself ignorant of the event.>

Fouquet began to walk about in his chamber with an uneasiness that became more and more serious.

<What do you decide upon, monseigneur?> said Gourville.

<If it were really as easy as you say, I would go to the king,> cried Fouquet. <But as I go to the Louvre, I will pass by the Hôtel de Ville. We shall see if the sentence is signed.>

<Incredulity! thou art the pest of all great minds,> said Gourville, shrugging his shoulders.

<Gourville!>

<Yes,> continued he, <and incredulity! thou ruinest, as contagion destroys the most robust health; that is to say, in an instant.>

<Let us go,> cried Fouquet; <desire the door to be opened, Gourville.>

<Be cautious,> said the latter, <the Abbé Fouquet is there.>

<Ah! my brother,> replied Fouquet, in a tone of annoyance; <he is there, is he? he knows all the ill news, then, and is rejoiced to bring it to me, as usual. The devil! if my brother is there, my affairs are bad, Gourville; why did you not tell me that sooner: I should have been the more readily convinced.>

<Monseigneur calumniates him,> said Gourville, laughing; <if he is come, it is not with a bad intention.>

<What, do you excuse him?> cried Fouquet; <a fellow without a heart, without ideas; a devourer of wealth.>

<He knows you are rich.>

<And would ruin me.>

<No, but he would have your purse. That is all.>

<Enough! enough! A hundred thousand crowns per month, during two years. \swear{Corbleu!} it is I that pay, Gourville, and I know my figures.> Gourville laughed in a silent, sly manner. <Yes, yes, you mean to say it is the king pays,> said the superintendent. <Ah, Gourville, that is a vile joke; this is not the place.>

<Monseigneur, do not be angry.>

<Well, then, send away the Abbé Fouquet; I have not a sou.> Gourville made a step towards the door. <He has been a month without seeing me,> continued Fouquet, <why could he not be two months?>

<Because he repents of living in bad company,> said Gourville, <and prefers you to all his bandits.>

<Thanks for the preference! You make a strange advocate, Gourville, to-day—the advocate of the Abbé Fouquet!>

<Eh! but everything and every man has a good side—their useful side, monseigneur.>

<The bandits whom the Abbé keeps in pay and drink have their useful side, have they? Prove that, if you please.>

<Let the circumstance arise, monseigneur, and you will be very glad to have these bandits under your hand.>

<You advise me, then, to be reconciled to the Abbé?> said Fouquet, ironically.

<I advise you, monseigneur, not to quarrel with a hundred or a hundred and twenty loose fellows, who, by putting their rapiers end to end, would form a cordon of steel capable of surrounding three thousand men.>

Fouquet darted a searching glance at Gourville, and passing before him,—<That is all very well; let M.~l'Abbé Fouquet be introduced,> said he to the footman. <You are right, Gourville.>

Two minutes after, the Abbé Fouquet appeared in the doorway, with profound reverence. He was a man of from forty to forty-five years of age, half churchman, half soldier,—a spadassin grafted upon an Abbé; upon seeing that he had not a sword by his side, you might be sure he had pistols. Fouquet saluted him more as elder brother than as a minister.

<What can I do to serve you, monsieur l'Abbé?> said he.

<Oh! oh! how coldly you speak to me, brother!>

<I speak like a man who is in a hurry, monsieur.>

The Abbé looked maliciously at Gourville, and anxiously at Fouquet, and said, <I have three hundred pistoles to pay to M.~de Bregi this evening. A play debt, a sacred debt.>

<What next?> said Fouquet bravely, for he comprehended that the Abbé Fouquet would not have disturbed him for such a want.

<A thousand to my butcher, who will supply no more meat.>

<Next?>

<Twelve hundred to my tailor,> continued the Abbé; <the fellow has made me take back seven suits of my people's, which compromises my liveries, and my mistress talks of replacing me by a farmer of the revenue, which would be a humiliation for the church.>

<What else?> said Fouquet.

<You will please to remark,> said the Abbé, humbly, <that I have asked nothing for myself.>

<That is delicate, monsieur,> replied Fouquet; <so, as you see, I wait.>

<And I ask nothing, oh! no,—it is not for want of need, though, I assure you.>

The minister reflected for a minute. <Twelve hundred pistoles to the tailor; that seems a great deal for clothes,> said he.

<I maintain a hundred men,> said the Abbé, proudly; <that is a charge, I believe.>

<Why a hundred men?> said Fouquet. <Are you a Richelieu or a Mazarin, to require a hundred men as a guard? What use do you make of these men?—speak.>

<And do you ask me that?> cried the Abbé Fouquet; <ah! how can you put such a question,—why I maintain a hundred men? Ah!>

<Why, yes, I do put that question to you. What have you to do with a hundred men?—answer.>

<Ingrate!> continued the Abbé, more and more affected.

<Explain yourself.>

<Why, monsieur the superintendent, I only want one valet de chambre, for my part, and even if I were alone, could help myself very well; but you, you who have so many enemies—a hundred men are not enough for me to defend you with. A hundred men!—you ought to have ten thousand. I maintain, then, these men in order that in public places, in assemblies, no voice may be raised against you; and without them, monsieur, you would be loaded with imprecations, you would be torn to pieces, you would not last a week; no, not a week, do you understand?>

<Ah! I did not know you were my champion to such an extent, monsieur le Abbé.>

<You doubt it!> cried the Abbé. <Listen, then, to what happened, no longer ago than yesterday, in the Rue de la Hochette. A man was cheapening a fowl.>

<Well, how could that injure me, Abbé?>

<This way. The fowl was not fat. The purchaser refused to give eighteen sous for it, saying that he could not afford eighteen sous for the skin of a fowl from which M.~Fouquet had sucked all the fat.>

<Go on.>

<The joke caused a deal of laughter,> continued the Abbé; <laughter at your expense, death to the devils! and the canaille were delighted. The joker added, <Give me a fowl fed by M.~Colbert, if you like! and I will pay all you ask.> And immediately there was a clapping of hands. A frightful scandal! you understand; a scandal which forces a brother to hide his face.>

Fouquet coloured. <And you veiled it?> said the superintendent.

<No, for so it happened I had one of my men in the crowd; a new recruit from the provinces, one M.~Menneville, whom I like very much. He made his way through the press, saying to the joker: <\swear{Mille barbes!} Monsieur the false joker, here's a thrust for Colbert!> <And one for Fouquet,> replied the joker. Upon which they drew in front of the cook's shop, with a hedge of the curious round them, and five hundred as curious at the windows.>

<Well?> said Fouquet.

<Well, monsieur, my Menneville spitted the joker, to the great astonishment of the spectators, and said to the cook:—<Take this goose, my friend, for it is fatter than your fowl.> That is the way, monsieur,> ended the Abbé, triumphantly, <in which I spend my revenues; I maintain the honour of the family, monsieur.> Fouquet hung his head. <And I have a hundred as good as he,> continued the Abbé.

<Very well,> said Fouquet, <give the account to Gourville, and remain here this evening.>

<Shall we have supper?>

<Yes, there will be supper.>

<But the chest is closed.>

<Gourville will open it for you. Leave us, monsieur l'Abbé, leave us.>

<Then we are friends?> said the Abbé, with a bow.

<Oh, yes, friends. Come, Gourville.>

<Are you going out? You will not stay to supper, then?>

<I shall be back in an hour; rest easy, Abbé.> Then aside to Gourville,—<Let them put to my English horses,> said he, <and direct the coachman to stop at the Hôtel de Ville de Paris.> 